\section{History.lean}

\code{History.lean} defines linear patch histories. A history is a list of
steps, where each step has a base commit id, a new commit id, and a structured
\code{Diff}. The file also defines shared-prefix splitting for two histories.

\subsection{Imports and Namespace}

This module depends on \code{Defns.lean} for \code{Diff} and
\code{Diff.metadataWellFormed}, and it stays in the shared
\code{VersionControl} namespace.

\subsection{Commit Ids and Patch Steps}

\begin{definition}[CommitId]
\code{[CommitId]} is an abbreviation for \code{String}. A commit id is just a
short string name such as \code{init}, \code{a1}, or \code{left1}.
\end{definition}

\begin{definition}[PatchStep]
\code{[PatchStep]} is a structure with fields \code{baseCommit},
\code{newCommit}, and \code{diff}. It represents one linear step from one
commit id to the next.
\end{definition}

\begin{definition}[PushHistory]
\code{[PushHistory]} is an abbreviation for \code{List PatchStep}. It
represents a linear sequence of patch steps.
\end{definition}

\begin{definition}[PatchStep.wellFormed]
\code{[PatchStep.wellFormed]} states that both commit ids are nonempty and that
the diff metadata is well formed.
\end{definition}

\subsection{Commit Paths and Distinctness}

\begin{definition}[commitPath]
\code{[commitPath]} converts a history into the sequence of commit ids it
visits. For \code{[init -> a1, a1 -> a2]}, the commit path is
\code{[init, a1, a2]}.
\end{definition}

\begin{definition}[distinctCommitIds]
\code{[distinctCommitIds]} states that the commit path has no duplicate commit
ids.
\end{definition}

\subsection{History Validity}

\begin{definition}[validHistory]
\code{[validHistory]} has type \code{PushHistory -> Prop}. A valid history must:

\begin{enumerate}
  \item use nonempty commit ids,
  \item chain adjacent steps correctly,
  \item keep all steps on the same file, and
  \item keep commit ids distinct across the whole linear history.
\end{enumerate}

This means histories with repeated commit ids are rejected rather than treated
as ambiguous ancestry.
\end{definition}

\subsection{Shared Prefix Logic}

\begin{definition}[commonPrefix]
\code{[commonPrefix]} returns the longest shared prefix of two lists.
\end{definition}

\begin{definition}[sharedCommitPath]
\code{[sharedCommitPath]} applies \code{commonPrefix} to the two histories'
\code{commitPath} values.
\end{definition}

\begin{definition}[sharedStepCount]
\code{[sharedStepCount]} returns the number of shared patch steps, computed from
the shared commit path length.
\end{definition}

\begin{definition}[CommonParentSplit]
\code{[CommonParentSplit]} packages the shared commit ids, the nearest common
ancestor commit id, and the left and right residual histories.
\end{definition}

\begin{definition}[splitAtNearestCommonParent?]
\code{[splitAtNearestCommonParent?]} computes the shared commit path. If there
is no shared commit id, it returns \code{none}. Otherwise it returns the nearest
common ancestor plus the residual histories obtained by dropping the shared
steps.
\end{definition}

\subsection{Theorems}

\begin{theorem}[commitPath\_nil]
\code{[commitPath_nil]} states that the commit path of the empty history is
\code{[]}.
\end{theorem}

\begin{theorem}[commitPath\_singleton]
\code{[commitPath_singleton]} states that the commit path of \code{[step]} is
\code{[step.baseCommit, step.newCommit]}.
\end{theorem}

\begin{theorem}[validHistory\_distinctCommitIds]
\code{[validHistory_distinctCommitIds]} states that every valid history has a
duplicate-free commit path.
\end{theorem}

\begin{theorem}[empty shared path gives no split]
If the shared commit path is empty, then there is no common-parent split.
\end{theorem}

\begin{theorem}[split\_eq\_some\_of\_sharedCommitPath]
\code{[split_eq_some_of_sharedCommitPath]} states the exact split returned when
the shared commit path is nonempty.
\end{theorem}

\begin{theorem}[validHistory\_tail]
\code{[validHistory_tail]} states that the tail of a valid nonempty history is
also valid.
\end{theorem}

\begin{theorem}[validHistory\_drop]
\code{[validHistory_drop]} states that dropping any number of steps from a valid
history leaves a valid history.
\end{theorem}

\begin{theorem}[validHistory\_leftResidual]
\code{[validHistory_leftResidual]} states that the left residual returned by a
valid split is valid.
\end{theorem}

\begin{theorem}[validHistory\_rightResidual]
\code{[validHistory_rightResidual]} states that the right residual returned by a
valid split is valid.
\end{theorem}

\subsection{Execution Pipeline}

The history pipeline is:

\begin{enumerate}
  \item represent a linear history as \code{PushHistory},
  \item check chaining, file consistency, and distinct commit ids with
    \code{validHistory},
  \item convert each history to a commit-id path with \code{commitPath},
  \item find the shared commit-id prefix with \code{sharedCommitPath},
  \item split at the nearest common ancestor with
    \code{splitAtNearestCommonParent?},
  \item pass the residual histories to \code{FullMerge.lean}.
\end{enumerate}
